\documentclass{article}
\usepackage[utf8]{inputenc}
\usepackage[T1]{fontenc}
\usepackage{graphicx}
\usepackage{algorithm}
\usepackage{algpseudocode}

\usepackage{listings}
\usepackage{xcolor}
\usepackage[margin=2cm]{geometry}
\usepackage{float}

\usepackage{amssymb} %Ajoute des symboles mathématiques comme R, Z, forall, exists, etc....
\usepackage{amsmath}
\usepackage{amsthm}

\usepackage{microtype} %utile pour gérer les espacements entre les caractères pour des fiches de cours

%\usepackage{titlesec} %Modifie les \section dans leur style (avec ce qu'il y a en bas de la ligne 21 à 25)
%\titleformat{\section}
 % {\normalfont\bfseries}
 % {(\thesection)}
  %{0.5em}
  %{}


\usepackage{enumitem} %permet d'énumerer avec un label en (a), (b), plutot que 1. 2. etc... 


\title{Révisions Exos Probabilités}
\author{Derrez Lorenzo}

\begin{document}
\maketitle

\section*{\textbf{Probas 1} (prépa)}

\subsection*{\textbf{CCINP 97} :}
\subsubsection*{Question 1 :}


\textbf{Loi de $X$ :} $X(\Omega) = \mathbb{N}$ et, pour tout $j \in \mathbb{N}$,
\begin{align*}
    \mathbf{P}(X=j) &= \sum_{k=0}^{+\infty} \mathbf{P}(X=j,\, Y=k) \\
                    &= \sum_{k \in \mathbb{N}} \frac{(j+k)\left(\frac{1}{2}\right)^{j+k}}{e\, j!\, k!} \\
                    &= \frac{1}{e\, j!\, 2^j} \left( j \sum_{k=0}^{+\infty} \frac{(1/2)^k}{k!} + \sum_{k=0}^{+\infty} \frac{k\,(1/2)^k}{k!} \right) \\
                    &= \frac{1}{e\, j!\, 2^j} \left( j\, e^{1/2} + \frac{1}{2}\, e^{1/2} \right) \\
                    &= \frac{1}{j!\, 2^j} \left( e^{-1/2} (j + \frac{1}{2}) \right)
\end{align*}

Par symétrie des rôles de $X$ et $Y$ :
$Y(\Omega) = \mathbb{N}$ et pour tout $k \in \mathbb{N}$ :
\begin{align*}
    \mathbf{P}(Y=k) &= \frac{1}{k!\, 2^k} \left( e^{-1/2} (k + \frac{1}{2}) \right)
\end{align*}

\subsubsection*{Question 2 :}
Remarquons,
\\
\begin{enumerate}
    \item $\mathbf{P}(X=0, Y=0) = 0$
    \item $\mathbf{P}(X=0)\mathbf{P}(Y=0) = \frac{e^{-1}}{4} $
\end{enumerate}

Par conséquent, les variables aléatoires $X$ et $Y$ ne sont pas indépendantes.


\end{document}